\chapter{Conclusiones}
\label{cap:conclusiones}

El presente Trabajo de Fin de Grado ha afrontado un reto que va más allá de la implementación técnica: demostrar que la domótica inteligente puede y debe ser un ámbito soberano para el usuario, sin dependencias de plataformas externas ni exposición de sus datos.

A lo largo del proceso de desarrollo, el proyecto ha madurado de forma significativa. La fase inicial, apoyada en la infraestructura universitaria para validar el enfoque agéntico, fue imprescindible para sentar las bases técnicas y arquitectónicas del sistema. La decisión de migrar posteriormente a una solución 100\% local no fue un retroceso, sino la materialización coherente de los principios que motivaron el proyecto desde su origen.

\section{Logros Alcanzados}

A fecha de redacción de este capítulo, el ecosistema Luxe Core AI cuenta con los siguientes hitos de implementación verificados sobre hardware real:

\begin{itemize}
    \item \textbf{Arquitectura de tres capas escalable:} el orquestador Node.js (OpenClaw npm) coordina los agentes de IA, los canales de usuario y la ejecución de \textit{drivers} Python de dispositivos, cada capa con interfaces bien definidas y gestionables de forma aislada.

    \item \textbf{Control local de dispositivos Tuya (RF-03 verificado):} el enchufe inteligente responde a comandos de activación y desactivación en menos de 2 segundos a través de la red air-gapped, utilizando el protocolo 3.5 de \texttt{tinytuya} sin pasar por los servidores de Tuya. La \textit{local key} se almacena exclusivamente en local. El procedimiento de \textit{Hotspot Spoofing} para el \textit{onboarding} de nuevos dispositivos está documentado y es reproducible.

    \item \textbf{Lectura de sensores ambientales BLE (RF-02 verificado):} los sensores Xiaomi LYWSD03MMC entregan temperatura, humedad y voltaje de batería mediante conexión GATT directa desde el servidor, sin necesidad de flashear firmware alternativo ni de obtener \textit{bind keys} de la nube de Xiaomi. El protocolo de lectura, basado en la característica \texttt{ebe0ccc1} del servicio propietario de Xiaomi, ha sido completamente documentado e implementado en el módulo \texttt{ble\_sensors.py}.

    \item \textbf{Control completo del ventilador FanLamp F8808 (RF-03 extendido):} se ha completado la ingeniería inversa del protocolo propietario Tencent, descubriendo que el ventilador utiliza anuncios BLE (y no conexiones GATT) para recibir comandos. El protocolo ha sido completamente documentado e implementado en \texttt{fanlamp\_bt.py}, un controlador Python que envía comandos al ventilador mediante \texttt{btmgmt} desde el adaptador Bluetooth del ordenador, sin necesidad de la aplicación del fabricante ni de hardware auxiliar. Se han mapeado y verificado 11 comandos (apagado general, ventilador on/off y 5 velocidades, luz on/off y modo noche). El coste económico de la integración ha sido de 0~\euro{}, utilizando exclusivamente las herramientas ya disponibles en el proyecto.

    \item \textbf{Orquestador OpenClaw desplegado:} el \textit{gateway} Node.js (OpenClaw npm v2026.5.27) se ejecuta como servicio \texttt{systemd} sobre el puerto 18789, accesible desde toda la subred \texttt{192.168.1.0/24}. Los agentes de IA utilizan Ollama local con los modelos \texttt{qwen2.5-coder:1.5b} (clasificador) y \texttt{qwen2.5-coder:7b-instruct} (razonador).

    \item \textbf{Canales de usuario operativos:} el \textit{gateway} de OpenClaw proporciona integración nativa con Telegram (BotFather API) y un \textit{dashboard} web accesible en \texttt{http://192.168.1.50:18789}, sin necesidad de procesos Python auxiliares.

    \item \textbf{Infraestructura de red aislada validada:} el router TP-Link autónomo con subred \texttt{192.168.1.0/24} y sin salida a Internet garantiza que ningún paquete generado por los dispositivos comerciales abandone el perímetro local. Esta arquitectura ha sido verificada operativamente durante todas las pruebas de integración.

    \item \textbf{Sistema de Zero-Inference masivo (462 patrones):} el Tier~0 del Model Router v2 se expandió de 127 a 462 patrones, cubriendo aproximadamente el 92\% de los comandos diarios con latencia inferior a 1~ms en RAM. Se incluyeron variantes con muletillas, dialectalismos andaluces y formas abreviadas, además de la cobertura completa de iluminación, ventilador, escenas y consultas.

    \item \textbf{Sistema de invalidación de caché:} se implementó un mecanismo de invalidación de la caché LRU del router que permite forzar el reanálisis de comandos cuando se actualizan los patrones o la configuración del sistema, sin necesidad de reiniciar el servicio.

    \item \textbf{Filtro de \texttt{log} de arranque ESP32:} se desarrolló un script que filtra y elimina los mensajes de \texttt{log} de arranque del ESP32 (\textit{boot log}), permitiendo una comunicación más limpia y fiable entre el servidor y el microcontrolador.

    \item \textbf{Disponibilidad del sistema:} durante el período de pruebas continuas, los servicios críticos (model-router, sensor-daemon y openclaw-gateway) permanecieron operativos sin experimentar caídas que requiriesen intervención manual. El sistema se mantuvo en funcionamiento ininterrumpido durante más de 24~horas en el servidor local, confirmando la estabilidad de la arquitectura y la eficacia de los mecanismos de auto-recuperación ante fallos transitorios.
\end{itemize}

\section{Estado Actual del Desarrollo}

\begin{table}[h]
\centering
\begin{tabular}{|l|c|p{6cm}|}
\hline
\textbf{Módulo} & \textbf{Estado} & \textbf{Observaciones} \\
\hline
Tuya Smart Plug & Completado & Validado en red air-gapped. ON/OFF alterna. \\
BLE Sensors (Xiaomi) & Completado & Lee 25°C / 59\% / 3048mV por GATT. \\
FanLamp Pro F8808 & Completado & Controlado vía BLE advertisements con 11 comandos verificados. \\
OpenClaw Orchestrator & Desplegado & \textit{Gateway} Node.js en puerto 18789 (LAN). Canales Telegram + web operativos. \\
Model Router v2 & Desplegado & Tres niveles: Tier~0 (462 patrones, \textasciitilde{}92\% cobertura), Tier~1 (1.5B clasificador), Tier~2 (7B razonador). Caché LRU con invalidación. \\
ESP32 HVAC Bridge & Pendiente & Firmware C++20 no implementado aún. \\
\hline
\end{tabular}
\caption{Estado de los módulos del ecosistema a la fecha de esta memoria.}
\end{table}

\section{Líneas de Trabajo Futuras}

\begin{itemize}
    \item \textbf{Desarrollo del firmware HVAC (ESP32 + LIN):} constituye la prioridad inmediata. El hardware está adquirido y cableado, y el repositorio \texttt{esphome-fujitsu-halcyon} de Omniflux~\cite{omniflux_esphome_fujitsu} proporciona una base sólida que nos ahorrará meses de ingeniería inversa del protocolo LIN. Durante el TFG no se completó por una combinación de alcance y logística doméstica: integrar y validar el firmware requiere poner la máquina fuera de servicio durante días, algo complicado en junio cuando el termómetro en Córdoba supera los 38~$^\circ$C y el aire acondicionado es más necesario que nunca. Queda como primer objetivo tras la entrega de la memoria.

    \item \textbf{Integración completa del bucle de control con OpenClaw:} conectar los módulos de sensores BLE, actuadores Tuya y control del ventilador FanLamp al bucle principal del orquestador para que el sistema tome decisiones de confort de forma autónoma en tiempo real.

    \item \textbf{Aprendizaje adaptativo:} incorporar modelos de \textit{reinforcement learning} para que el sistema aprenda los hábitos del usuario y anticipe sus preferencias de confort, siguiendo la línea iniciada con el autoaprendizaje del Smart Advisor.

    \item \textbf{Expansión sensorial:} integrar un sensor de movimiento por infrarrojos o radar (como el LD2410) para mejorar la detección de presencia sin depender del teléfono móvil. También se plantea incorporar un dispositivo con micrófono y altavoz que permita la comunicación por voz con el sistema, eliminando la necesidad de escribir comandos.

    \item \textbf{Base de datos vectorial:} cuando el ecosistema escale a decenas de dispositivos y cientos de comandos, incorporar una base de datos vectorial (Qdrant o ChromaDB) para mejorar la recuperación semántica.

    \item \textbf{Expansión del \textit{dashboard}:} extender el \textit{dashboard} nativo de OpenClaw con visualización de analíticas energéticas e histórico de confort ASHRAE 55~\cite{ashrae55_2023}.

    \item \textbf{Expansión de hardware:} integrar nuevos protocolos como Zigbee o Z-Wave para ampliar el catálogo de dispositivos compatibles.

    \item \textbf{Automatización de infraestructura:} containerizar los servicios del ecosistema (OpenClaw, Model Router, Sensor Daemon) mediante Docker y orquestarlos con \texttt{docker-compose}, facilitando el despliegue en otros hogares. Incorporar aprovisionamiento automatizado con Ansible y un \textit{pipeline} CI/CD que ejecute la suite de tests de integración antes de cada despliegue, siguiendo prácticas de IaC (\textit{Infrastructure as Code}).
\end{itemize}

\section{Reflexión Personal}

Más allá de los logros técnicos, este proyecto ha sido un proceso de aprendizaje continuo. Cada obstáculo (desde un puerto Ethernet mal conectado hasta un protocolo BLE sin documentación) se convirtió en una oportunidad para aprender algo nuevo. He disfrutado especialmente la fase de ingeniería inversa del ventilador FanLamp, donde la combinación de análisis de código, experimentación con hardware y paciencia permitió descifrar un protocolo que el fabricante no había previsto que nadie entendiera.

La decisión de mantener todo el sistema en local, sin depender de la nube, no fue solo técnica: fue también una declaración de principios. La tecnología del hogar no debería exigir ceder la privacidad a cambio de comodidad. Este TFG demuestra que es posible tener ambas cosas.
